\chapter{Conclusions}
\label{ch:conclusions}

\section{Summary}
\label{sec:conclusions_summary}

This thesis addressed two interconnected problems in biometric template protection:
(1)~the AND-fusion gating penalty that causes every existing multimodal fuzzy vault
to perform worse than unimodal fingerprint recognition, and (2)~the single-point
of compromise inherent in centralized biometric databases. The proposed solutions---
the \textbf{Iris-Boosted Fuzzy Vault (IBFV)} and its decentralized extension
\textbf{D-IBFV}---resolve both problems simultaneously, with formal proofs and
extensive experimental validation.

The four key results of this work are:

\begin{enumerate}
  \item \textbf{Zero-penalty multimodal fusion}: IBFV's conditional iris-boosting
  architecture guarantees that the multimodal system never performs worse than
  unimodal fingerprint recognition, regardless of iris quality. This was verified
  empirically across all 144 Setup 2 configurations (zero violations) and proved
  formally in Theorem~1 (Appendix~A).

  \item \textbf{Significant EER reduction}: When iris recovery succeeds ($\sim$62\%
  of genuine attempts), IBFV reduces EER by up to \textbf{61.3\%} relative to
  the unimodal baseline (FVC2002~DB3). In Setup~1, 113 of 144 configurations
  improve or tie; 31 configurations on the easy databases (DB1/DB2) show marginal
  elevation ($\leq 0.07$~pp) attributable to the non-zero iris impostor FAR
  ($\approx 0.10\%$). The improvement is largest on challenging databases where the
  fingerprint alone struggles at the margin.

  \item \textbf{Decentralized template protection with D-IBFV}: Applying
  $(t,n)$-Shamir Secret Sharing to the IBFV vault, with shares stored on IPFS
  and CIDs anchored on Hyperledger Fabric and EVM-compatible blockchains,
  eliminates single-server trust. Across 30{,}000 reconstruction trials, the
  reconstructed vault is bit-for-bit identical to the original. The system tolerates
  up to $n{-}t$ node failures and provides information-theoretic security against
  compromise of any $t{-}1$ shares (Theorem~2).

  \item \textbf{Practical performance}: D-IBFV verification completes in
  $\sim$56--70\,ms end-to-end (including blockchain CID lookup and IPFS share
  retrieval), with the biometric verification taking only $\sim$12\,ms. The system
  is compatible with ISO/IEC~24745 irreversibility, unlinkability ($D_\mathrm{sys}
  \leq 0.03$), and revocability requirements.
\end{enumerate}

\section{Future Work}
\label{sec:future_work}

\begin{enumerate}
  \item \textbf{Adaptive $K$ selection}: The number of bonus points $K$ is
  currently fixed at enrollment. An adaptive scheme that adjusts $K$ based on
  real-time iris quality scores (e.g., BRISQUE, pupil dilation, coverage fraction)
  could further optimize the boost while minimizing any latency overhead.

  \item \textbf{Co-acquired multimodal datasets}: All experiments used chimeric
  (artificially paired) fingerprint-iris data. Evaluation on genuinely co-acquired
  multimodal databases (e.g., a custom collection at NIT Warangal or the
  WVU multimodal database) would validate performance under correlated intra-subject
  variation.

  \item \textbf{Production blockchain deployment}: Deploying D-IBFV on a public
  Layer-2 network (Polygon PoS, Arbitrum, StarkNet) would provide real-world
  latency and cost measurements under network congestion, mempool competition,
  and variable gas prices---conditions not simulated by Ganache.

  \item \textbf{Verifiable Secret Sharing}: Replacing Shamir's basic scheme with
  a Verifiable Secret Sharing (VSS) protocol \cite{pedersen1991non} would allow
  each shareholder to prove that their share is correct without reconstructing
  the vault. This would enable malicious node detection and is a natural
  extension for adversarial deployment environments.

  \item \textbf{Additional modalities}: The IBFV architecture is not specific to
  iris and fingerprint. Any biometric capable of producing a stable binary key
  (face \cite{hao2006combining}, finger vein, palm print) can serve as the booster
  modality. A systematic study of which modality combinations maximize the boost
  is a natural extension.

  \item \textbf{On-device enrollment}: Migrating the enrollment computation to a
  trusted execution environment (ARM TrustZone, Intel TDX) would allow the vault
  to be constructed on the user's device, with only shares uploaded, further
  reducing the attack surface.
\end{enumerate}


